% Part: first-order-logic
% Chapter: completeness
% Section: downward-ls

\documentclass[../../../include/open-logic-section]{subfiles}

\begin{document}

\olfileid{fol}{com}{dls}
\olsection{De stelling van L\"owenheim--Skolem}

De stelling van L\"owenheim--Skolem zegt dat als een theorie een oneindig
model heeft, zij tevens een model heeft dat eindig of !!{denumerable} is.
Hieruit volgt onmiddellijk dat de eerste-ordelogica niet kan uitdrukken
dat de grootte van !!a{structure} !!{nonenumerable} is: elke !!{sentence}
of verzameling !!{sentence}s die in alle !!{nonenumerable} !!{structure}s
wordt vervuld, wordt ook in een !!{enumerable} structuur vervuld.

\begin{thm} 
\ollabel{thm:downward-ls} Als $\Gamma$ consistent is, heeft zij
!!a{enumerable} model; zij is dus vervulbaar in !!a{structure}
waarvan het domein eindig of !!{denumerable} is.
\end{thm}

\begin{proof}
Als $\Gamma$ consistent is, heeft de !!{structure}~$\Struct M$ die door
het bewijs van de volledigheidsstelling wordt geleverd een domein
$\Domain{M}$ dat niet groter is dan de verzameling termen van de
taal~$\Lang L$. De structuur $\Struct M$ is dus hoogstens
!!{denumerable}.
\end{proof}

\begin{thm}
\ollabel{noidentity-ls} Als $\Gamma$ een consistente verzameling
!!{sentence}s is in de taal van de eerste-ordelogica zonder identiteit,
dan heeft zij !!a{denumerable} model; zij is dus vervulbaar in
!!a{structure} waarvan het domein oneindig en !!{enumerable} is.
\end{thm}

\begin{proof}
Als $\Gamma$ consistent is en geen zinnen bevat waarin identiteit
voorkomt, heeft de !!{structure}~$\Struct M$ die door het bewijs van de
volledigheidsstelling wordt geleverd een domein $\Domain{M}$ dat gelijk
is aan de verzameling termen van de taal~$\Lang L'$. De structuur
$\Struct{M}$ is dus !!{denumerable}, aangezien $\Trm[L']$ dat is.
\end{proof}

\begin{ex}[De paradox van Skolem]
De verzamelingenleer van Zermelo--Fraenkel~$\Log{ZFC}$ is een zeer
krachtig raamwerk waarin vrijwel alle wiskundige uitspraken kunnen
worden uitgedrukt, waaronder uitspraken over de grootte van
verzamelingen. Zo kan $\Log{ZFC}$ bijvoorbeeld bewijzen dat de
verzameling~$\Real$ van reële getallen !!{nonenumerable} is, en bewijst
de theorie de stelling van Cantor dat de machtsverzameling van elke
verzameling groter is dan de verzameling zelf. Als $\Log{ZFC}$
consistent is, zijn al haar modellen oneindig. Bovendien bevatten zij
allemaal !!{element}s waarvan de theorie zegt dat zij !!{nonenumerable}
zijn, zoals het element dat de stelling van~$\Log{ZFC}$ waar maakt dat de
machtsverzameling van de natuurlijke getallen bestaat. Volgens de
stelling van L\"owenheim--Skolem heeft $\Log{ZFC}$ ook !!{enumerable}
modellen: modellen die ``!!{nonenumerable}'' verzamelingen bevatten maar
zelf !!{enumerable} zijn.
\end{ex}

\end{document}
